% page 391 of the facsimile (image p-390 of the scan)
% block 175.3 x 242.0 mm ; left margin 26.7 mm
\pgc{5.080mm}{0.000mm}{
\l{\cel{58}383}
\l{\cel{44}17\cel{2}19\cel{3}3}
\l{\cel{4}\sou{morfino}. (\sou{kemio)}. Alkaloido di opiumo, C\cel{3}N . NO\cel{2}+H\cel{2}O.;}
\l{\cel{7}narkotigivo tre efikiva, uzata medicinale por mingravigar}
\l{\cel{7}la sufri;drogo danjeroza pro kustumesko. - DEFIRS.}
\l{}
\l{\cel{4}morfogena. (\sou{biol.}) Dicesas pri la efiki e la agenti qui influas}
\l{\cel{7}la formo e la strukturo di la organismi. - DEFIS.}
\l{}
\l{\cel{4}morfogenezo. (biol.) Cienco di la kauzi qui determinas la formo}
\l{\cel{7}e la strukturo di la organismi, e di la +leyi (legi) qui}
\l{\cel{7}rezumas li. - DEFIS.}
\l{}
\l{\cel{4}morfologio. I. (\sou{filoz.}) Cienco di la formi diversa di materio.}
\l{\cel{7}- II. (\sou{gram.}) Cienco di la formi diversa gramatikala di}
\l{\cel{7}la vorti. - DEFIRS.}
\l{}
\l{\cel{4}\sou{morfoplasmo}. (\sou{biol}.)Un ek la substanci protoplasma di la}
\l{\cel{7}celulo-korpo di la ovulo, qua efikas a la karakterizivi di}
\l{\cel{7}la genitori. -}
\l{}
\l{\cel{4}\sou{morfoso}. (\sou{biol.}) Reflexo elementala, ne-nervala, qua esas res-}
\l{\cel{7}pondo rapida ad ecito simpla di la substanco vivanta, e qua}
\l{\cel{7}su manifestas per ula formo od equilibro-posturo di ta sub-}
\l{\cel{7}stanco vivanta en la medio cirkondanta. - DEFIS.}
\l{}
\l{\cel{4}\sou{morge}. En la dio qua sequas nemediate olta pri qua onu parolas.}
\l{\cel{7}- DE.}
\l{}
\l{\cel{4}\sou{morganatika}. (\sou{yurocienco}).(Pariajo) quan princo o sinioro paktas}
\l{\cel{7}kun persono di rango inferiora e qua ne havas omna efekti}
\l{\cel{7}civila. - DEFIRS.}
\l{}
\l{\cel{4}\sou{morilio}. Fero-peco, di qua un extremajo kapo-forma esas}
\l{\cel{7}ornita, e quan onu pozas ye singla latero di la herdo di}
\l{\cel{7}kameno, por subtenar la kombusto-ligno.\cel{2}- S.}
\l{}
\l{\cel{4}\sou{morso}. Levero muntita en la boko di la kavalo, e qua, per efikar}
\l{\cel{7}a la stangi, uzesas por direktar lu. - FI.}
\l{}
\l{\cel{4}\sou{mortar}. (\sou{netrans}.) Cesar vivar. (\sou{anke metaf.}). - DEFIS.}
\l{}
\l{\cel{4}\sou{mortadelo.}\cel{2}Grosa socisego ek hepati e sen-grasa karni (pistita).}
\l{\cel{7}- DFIS.}
\l{}
\l{\cel{4}mortero. Sablo o cimento mixita a kalko, e dilutita en aquo,}
\l{\cel{7}quan onu uzas por kunjuntar la petri e quadri di konstrukturo,}
\l{\cel{7}di edifico. - DEFIS.}
\l{}
\l{\cel{4}mortezo.\cel{2}(\sou{teknol}.) Entaliuro en peco de ligno o metalo, por}
\l{\cel{7}recevar tenono facita en altra peco, tale ke formacesas}
\l{\cel{7}asembluro. - EFS.}
\l{}
\l{\cel{4}mortifikar. (\sou{trans.}) Igar (ulu) mortar ye su ipsa, per domtar}
\l{\cel{7}sua impulsi karnala, e per humiligar la spirito per asket-}
\l{\cel{7}alaji. - EFIRS.}
}
